% from main.tex
%%%%%%%%%%  scan idx 548  |  folio 535  %%%%%%%%%%
Soit $\beta_y:A\to K_\bullet(X)$ l'homomorphisme ainsi défini. Si l'on désigne par $I$ l'idéal $\Filt^1(X)/\Filt^{j+1}(X)$ de $A$, on a $I^{j+1}=0$. Gardant les notations du lemme 2.1.4 on trouve un homomorphisme de groupes abéliens
\[
\gamma:\Pic(X)\longrightarrow P
\]
qui associe à la classe de tout $\mathcal O_X$-module inversible $L$ dans $\Pic(X)$ l'élément $j(\cl^\circ(L))$ de $P$.

On tire immédiatement des lemmes 2.1.2 et 2.1.4 :

\textbf{Proposition 2.2.1.} Pour tout schéma noethérien $X$ et pour tout $y\in\Filt_j(X)$ ($j\in\mathbb Z$) l'application composée
\[
\Pic(X)\longrightarrow K_\bullet(X)\otimes_{\mathbb Z}\mathbb Q,
\]
définie par
\[
\cl_P(L)\rightsquigarrow (\cl^\circ(L)y)\otimes1
\]
est une application polynôme de degré $\leq j$.

\subsection*{2.3}
Si $X$ est un schéma propre sur un corps $k$ et si $g:X\to\Spec(k)$ est le morphisme structural, alors l'homomorphisme
\[
g_*:K_\bullet(X)\longrightarrow K_\bullet(\Spec(k))\cong\mathbb Z
\]
associe à la classe $\cl_\bullet(F)$ de tout $\mathcal O_X$-module cohérent $F$ la caractéristique d'Euler-
